% Part: lambda-calculus
% Chapter: introduction
% Section: currying

\documentclass[../../../include/open-logic-section]{subfiles}

\begin{document}

\olfileid{lam}{int}{cur}
\olsection{Currying}

Abstraksi-$\lambd$~$\lambd[x][M]$ makili fungsi
siji argumen. Iki dadi watesan nalika kita arep
netepake fungsi kang nampa luwih saka siji argumen.
Salah siji cara yaiku nggedhekake kalkulus-$\lambd$
supaya bisa mbentuk pasangan, triplèt, lan sateruse;
yen mangkono, umpamane, fungsi telung argumen
$\lambd[x][M]$ bakal nampa sawijining triplèt
minangka argumene. Nanging luwih praktis yen
kita nggunakake \emph{Currying}.

Gatekna sawijining tuladha. Sedhela wae, upamakna
ana operasi~$+$ ing kalkulus-$\lambd$.
Fungsi tambah iku fungsi $2$ argumen, yaiku
nampa rong argumen. Nanging abstraksi-$\lambd$
mung ngasilake fungsi siji argumen: sintaksise
ora ngidini ekspresi kaya
$\lambd[(x,y)][(x+y)]$. Senadyan mangkono,
kita bisa nimbang fungsi siji argumen~$f_x(y)$
kang diwenehake dening $\lambd[y][(x+y)]$,
yaiku nambahake $x$ marang argumen tunggale~$y$.
Satemene iki dudu siji fungsi, nanging kulawarga
fungsi kang beda-beda, ``tambah~$x$'',
siji kanggo saben wilangan~$x$. Saiki kita bisa
netepake fungsi siji argumen liyane~$g$ minangka
$\lambd[x][f_x]$. Yen diaplikasikake marang
argumen $x$, $g(x)$ ngasilake fungsi~$f_x$;
dadi nilai-nilaine uga fungsi. Saiki yen
kita ngaplikasikake $g$ marang $x$, banjur
ngaplikasikake asile marang~$y$, kita oleh
$(g(x))y = f_x(y) = x+y$. Kanthi cara iki,
fungsi siji argumen~$g$ bisa nindakake pakaryan
kang padha karo fungsi tambah rong argumen.
``Currying'' mung jeneng kanggo cara
ngowahi fungsi rong argumen dadi fungsi siji
argumen kang nilai-nilaine fungsi siji argumen.

Ing ngisor iki tuladha kang trep ing sintaksis
kalkulus-$\lambd$. Kepiye carane makili
fungsi $f(x,y) = x$? Yen kita arep netepake
fungsi kang nampa rong argumen lan ngasilake
argumen kapisan, kita bisa nulis
$\lambd[x][\lambd[y][x]]$, kang teges
harfiahe fungsi kang nampa argumen~$x$
lan ngasilake fungsi~$\lambd[y][x]$.
Fungsi $\lambd[y][x]$ nampa argumen
liyane~$y$, nanging ora migunakake argumen
iku lan tansah ngasilake~$x$.
Gatekna apa kang kedadeyan nalika
$\lambd[x][\lambd[y][x]]$
diaplikasikake marang rong argumen:
\begin{align*}
  (\lambd[x][\lambd[y][x]])MN
  \bredone & (\lambd[y][M])N \\
  \bredone & M
\end{align*}

Umume, kanggo nulis fungsi kanthi parameter
$x_1$, \dots,~$x_n$ kang ditetepake
dening sawijining suku~$N$, kita bisa nulis
$\lambd[x_1][\lambd[x_2][\ldots\lambd[x_n][N]]]$.
Yen kita ngaplikasikake $n$ argumen
marang fungsi iki, kita oleh:
\begin{multline*}
  (\lambd[x_1][\lambd[x_2][\ldots\lambd[x_n][N]]]) M_1 \dots M_n \\
  \begin{aligned}
  \bredone {} & (\Subst{(\lambd[x_2][\ldots\lambd[x_n][N]])}{M_1}{x_1}) M_2
  \dots M_n\\
   \eqs {} & (\lambd[x_2][\ldots\lambd[x_n][\Subst{N}{M_1}{x_1}]]) M_2
            \dots M_n \\
  \vdots & \\
  \bredone {} &\Subst{\Subst{N}{M_1}{x_1}\ldots}{M_n}{x_n}
  \end{aligned}
\end{multline*}
Larik pungkasan teges harfiahe nyelehake
$M_i$ minangka pangganti $x_i$ ing awak
definisi fungsi, persis kaya kang dikarepake
nalika fungsi diaplikasikake marang luwih saka
siji argumen.
\end{document}
